% Paper 2, §5 — Distributed training
% Anchors: kb/notes/training/distributed-training.md
%          kb/excerpts/{torchtitan2024,deepseek-v3-training,muon-moonlight2025}.md

\section{Distributed training}
\label{sec:distributed-training}

A frontier-scale pre-training run no longer chooses ``a parallelism.''
It picks five sizes simultaneously:
\begin{equation}
G \;=\; d \cdot t \cdot c \cdot p \cdot e,
\label{eq:5d-mesh}
\end{equation}
where $G$ is the total worker count and $(d, t, c, p, e)$ are the
degrees of \emph{\glsterm{data-parallel}{Data Parallel}} (\glsterm{data-parallel}{DP}), \emph{\glsterm{tensor-parallel}{Tensor Parallel}} (\glsterm{tensor-parallel}{TP}, often
fused with \glsterm{sequence-parallel}{Sequence Parallel}), \emph{\glsterm{context-parallel}{Context Parallel}} (\glsterm{context-parallel}{CP}),
\emph{\glsterm{pipeline-parallel}{Pipeline Parallel}} (\glsterm{pipeline-parallel}{PP}), and \emph{\glsterm{expert-parallel}{Expert Parallel}} (\glsterm{expert-parallel}{EP})
respectively \citep[\S 1]{torchtitan2024}.\footnote{KB excerpt:
\texttt{kb/excerpts/torchtitan2024.md\#sec-1-axes}.} Each axis is a
distinct collective with its own bandwidth-vs-latency profile, its own
memory savings, and its own interaction with model architecture. The
thesis of this section is that the 2024--2026 frontier has converged
on these five axes as a single composable surface --- a $5$D
\texttt{DeviceMesh} \citep[\S 2.1.1]{torchtitan2024}\footnote{KB
excerpt: \texttt{kb/excerpts/torchtitan2024.md\#sec-2-1-1-meta}.} on
which a parallelism plan is a tuple of axis sizes --- and that
designing a frontier training run is now picking that tuple under
joint constraints of memory, interconnect topology, and the model
graph. We work each axis through its per-step communication, per-layer
communication, per-microbatch activation, and in-flight activation
costs; collect a memory table; walk \glsterm{precision-pinning}{DeepSeek-V3}'s plan as the
load-bearing case study; and close on three live contradictions
(\glsterm{tensor-parallel}{TP}-yes-vs-\glsterm{tensor-parallel}{TP}-no, optimal-plan-from-the-graph, the Blackwell NVL72
inversion) that §14 picks back up.

\subsection{The five axes}
\label{sec:dt-axes}

\paragraph{DP / FSDP / ZeRO.} Pure \glsterm{data-parallel}{DP} replicates \glsterm{parameters}{parameters} $\theta
\in \R^P$, gradients, and Adam moments on each of $d$ workers, all-
reducing gradients once per step. The optimizer-state footprint at
FP32-master-plus-Adam is $4P$ words per worker (\glsterm{parameters}{parameters} + gradients
+ first moment + second moment) \citep[\S 2]{kalamkar2019-bfloat16}.
\glsterm{zero-1-zero-2-zero-3}{ZeRO-1 / ZeRO-2 / ZeRO-3} progressively shard optimizer state /
gradients / \glsterm{parameters}{parameters} across the $d$ workers, trading per-step
collective volume for memory \citep[\S 2.1.2]{torchtitan2024}.\footnote{KB
excerpt: \texttt{kb/excerpts/torchtitan2024.md\#sec-2-1-2-fsdp2}.}
\glsterm{dtensor-devicemesh}{PyTorch}'s FSDP1 implemented full ZeRO-3 via a ``\texttt{FlatParameter}''
representation that bucketed all of a layer's \glsterm{parameters}{parameters} into one
contiguous tensor; \glsterm{fsdp2}{FSDP2} replaces that with per-parameter
\texttt{DTensor} sharding, which composes cleanly with \glsterm{tensor-parallel}{TP}/\glsterm{context-parallel}{CP}/\glsterm{pipeline-parallel}{PP} and
delivers ``often $7\%$ lower per-GPU memory requirement versus FSDP1''
plus ``an average of $1.5\%$'' throughput gain \citep[\S
2.1.2]{torchtitan2024}. Per-step cost: one all-gather of \glsterm{parameters}{parameters}
($\Theta(P)$) and one reduce-scatter of gradients ($\Theta(P)$),
amortized across the model's microbatches.

\paragraph{TP + SP.} \glsterm{tensor-parallel}{Tensor Parallel} \citep[\S 2.1.3]{torchtitan2024}
splits each \texttt{Linear} layer's weight matrix $W \in \R^{d_{\text{in}}
\times d_{\text{out}}}$ into $t$ row- or column-shards. \glsterm{dtensor-devicemesh}{PyTorch}
expresses this as the \texttt{RowwiseParallel} / \texttt{ColwiseParallel}
DTensor placements: column-shard the QKV projection, row-shard the
attention output projection, and a forward all-reduce stitches the
shards back. One all-reduce per linear in forward and one per linear
in backward --- on the order of $2L$ per-layer collectives per step.
\glsterm{sequence-parallel}{Sequence Parallel} (\glsterm{sequence-parallel}{SP}) is the always-paired companion: normalization
and dropout layers are sharded along the token dimension instead of
replicated, since otherwise their activations are \glsterm{tensor-parallel}{TP}-replicated and
become memory-dominant on long context \citep[\S 2.1.3]{torchtitan2024}.
Per-GPU activation under \glsterm{tensor-parallel}{TP}+\glsterm{sequence-parallel}{SP} shrinks by roughly $1/t$ along the
sharded dimension. Naive \glsterm{tensor-parallel}{TP} serializes the matmul against its
collective; \emph{Async \glsterm{tensor-parallel}{TP}} fractionalizes the sharded matmul into
chunks and overlaps each chunk's compute with the previous chunk's
collective via an intra-node \texttt{SymmetricMemory} buffer that
permits direct P2P writes between GPUs in a node \citep[\S
2.2.3]{torchtitan2024}.\footnote{KB excerpt:
\texttt{kb/excerpts/torchtitan2024.md\#sec-2-2-3-async-tp}.}

\paragraph{CP (Context Parallel).} \glsterm{context-parallel}{CP} shards along the
\emph{token-position} axis: a $\mathsf{B}\!\times\!\mathsf{S}\!\times\!\mathsf{D}$
\glsterm{residual-stream}{residual stream} becomes $\mathsf{B}\!\times\!(\mathsf{S}/c)\!\times\!\mathsf{D}$
on each of $c$ workers. Attention is partitioned via a ring-attention-
style schedule: each worker holds its sequence chunk's queries and
streams keys/values around a ring of $c$ workers, accumulating partial
\glsterm{softmax}{softmax} denominators as it goes, so the full \glsterm{softmax}{softmax} is computed
without any single worker materializing the $\mathsf{S}\!\times\!\mathsf{S}$
attention matrix. The headline empirical result is from
\citet[\S 2.1.5]{torchtitan2024}\footnote{KB excerpt:
\texttt{kb/excerpts/torchtitan2024.md\#sec-2-1-5-cp}.}: ``on Llama 3.1
8B with 8 H100 GPUs, using \glsterm{context-parallel}{CP} enabled training at context lengths up
to $262{,}144$ tokens, achieving minor \glsterm{mfu}{MFU} degradation as \glsterm{context-parallel}{CP} degree
increases.'' Without \glsterm{context-parallel}{CP} the same setup OOMs.

\paragraph{PP.} \glsterm{pipeline-parallel}{PP} partitions the layer stack into $p$ stages, one
stage per worker group, with microbatches flowing stage-to-stage via
P2P send/recv \citep[\S 2.1.4]{torchtitan2024}. Per-GPU parameter
memory drops as $1/p$, but in-flight activations multiply by $p$
because the pipeline must hold one microbatch's worth of activation per
stage. The classic 1F1B schedule \citep[after][]{harlap2018-pipedream}
has a ``bubble'' on the first forward pass through all $p$ stages and
on the last backward pass, totaling $2(p-1)$ stage-times of idle. The
schedule lineage tightens this:
\begin{equation}
\text{1F1B:}\;\; (p-1)(F + B), \qquad
\text{ZeroBubble (ZB1P):}\;\; (p-1)(F + B - 2W), \qquad
\text{DualPipe:}\;\; \!\bigl(\tfrac{p}{2} - 1\bigr)(F\&B + B - 3W),
\label{eq:bubble-formulas}
\end{equation}
verbatim from \citet[\S 3.2.1, table]{deepseek-v3}\footnote{KB excerpt:
\texttt{kb/excerpts/\glsterm{precision-pinning}{deepseek-v3}-training.md\#sec-3-2-1-bubble}.}, where
$F$ is the forward microbatch time, $B$ the (combined) backward time,
and $W$ the weight-gradient sub-step that \glsterm{zerobubble}{ZeroBubble} and \glsterm{dualpipe}{DualPipe}
extract from $B$. \glsterm{zerobubble}{ZeroBubble} reorders the freshly-isolated $W$ into
the bubble; \glsterm{dualpipe}{DualPipe} additionally feeds the pipeline from both ends
simultaneously so the first-microbatch ramp on one side overlaps with
the last-microbatch ramp on the other (we walk \glsterm{dualpipe}{DualPipe} in detail in
\cref{sec:dt-deepseek}). \deepencite{ZeroBubble (Qi et al.\ 2023b)
formal $W$-decomposition --- KB has only the bubble formula transcribed
from \citet[\S 3.2.1]{deepseek-v3}.}

\paragraph{EP (Expert Parallel).} For an MoE layer with $E$ routed
experts, total expert \glsterm{parameters}{parameters} $E \cdot P_E$ are too large for a
single GPU at frontier scale --- \glsterm{precision-pinning}{DeepSeek-V3} has $256$ routed experts
per layer \citep[\S 3.2]{deepseek-v3}. \glsterm{expert-parallel}{EP} shards experts across $e$
workers; each token's forward pass becomes:
\begin{equation}
\underbrace{\text{all-to-all}_{\text{dispatch}}}_{\text{token}\to\text{expert host}}
\;\to\;
\underbrace{\text{expert FFN compute}}_{\text{local}}
\;\to\;
\underbrace{\text{all-to-all}_{\text{combine}}}_{\text{expert host}\to\text{token}}
\label{eq:ep-pattern}
\end{equation}
\citep[\S 3.2; \S 3.2.2]{deepseek-v3}. \glsterm{expert-parallel}{EP} is bandwidth-heavy and
typically cross-node; in fact, per
\citet[\S 3.2.1]{deepseek-v3}\footnote{KB excerpt:
\texttt{kb/excerpts/\glsterm{precision-pinning}{deepseek-v3}-training.md\#sec-3-2-1-dualpipe}.}, the
cross-node \glsterm{expert-parallel}{EP} all-to-all gives \glsterm{precision-pinning}{DeepSeek-V3} ``an inefficient
computation-to-communication ratio of approximately 1:1,'' which is
exactly what \glsterm{dualpipe}{DualPipe} is designed to hide.

\subsection{Per-axis cost model}
\label{sec:dt-cost-model}

Let $P$ denote total \glsterm{parameters}{parameters} (FP32-master + Adam moments give $4P$
optimizer-state words) and $A$ the peak per-microbatch activation
memory of one stage of the model. \cref{tab:dt-memory} summarizes the
per-GPU memory and communication cost under each axis combination.

\begin{table}[h]
\centering
\small
\begin{tabular}{lcccc}
\toprule
Axis combination & Per-GPU param & Per-GPU optim state & Per-GPU activation & Per-step / per-layer comm \\
\midrule
Pure \glsterm{data-parallel}{DP} (ZeRO-0)            & $P$     & $\sim 4P$           & $A$              & all-reduce $\Theta(P)$ / step \\
ZeRO-1 (optim shard)        & $P$     & $4P/d$              & $A$              & all-gather + reduce-scatter $\Theta(P)$ / step \\
ZeRO-2 (grad shard)         & $P$     & $4P/d$ + grads$/d$  & $A$              & + reduce-scatter $\Theta(P)$ / step \\
ZeRO-3 / \glsterm{fsdp}{FSDP}-full          & $P/d$   & $4P/d$              & $A$              & all-gather param shards / layer \\
\glsterm{tensor-parallel}{TP} only (size $t$)          & $P/t$   & $4P/t$              & $A/t$ (sharded)  & 2 all-reduces / linear ($\sim 2L$/step) \\
\glsterm{pipeline-parallel}{PP} only (stages $p$)        & $P/p$   & $4P/p$              & $A \cdot p$ (in-flight) & P2P send/recv $\Theta(\text{act})$ / boundary \\
\glsterm{context-parallel}{CP} (size $c$)               & $P$     & $4P$                & $A/c$            & ring KV pass $\Theta(\mathsf{S}\!\cdot\!\mathsf{D})$ / attention \\
\glsterm{expert-parallel}{EP} (size $e$)               & $P_{\text{shared}} + P_E E/e$ & matched & $\sim A$ + dispatch buffer & 2 all-to-alls / MoE layer \\
\bottomrule
\end{tabular}
\caption{Per-GPU memory and communication cost under each parallelism
axis. $P$ is parameters, $A$ is peak per-microbatch activation per
stage, $L$ is layer count, $\mathsf{S}$ is sequence length,
$\mathsf{D}$ is model dim. Optimizer state $4P$ assumes
FP32-master-plus-Adam-moments \citep[\S 2]{kalamkar2019-bfloat16};
\cref{sec:mixed-precision} discusses the BF16-Adam-moments and
fine-grained-FP8 reductions. After
\citet[\S 2.1]{torchtitan2024} and \citet[\S 3.2]{deepseek-v3}.}
\label{tab:dt-memory}
\end{table}

Two structural tensions live in \cref{tab:dt-memory}. The first is
\emph{\glsterm{pipeline-parallel}{PP} saves \glsterm{parameters}{parameters} but multiplies in-flight activations by
$p$}: the bookkeeping for $p$ microbatches' worth of activation per
stage is exactly why \glsterm{dualpipe}{DualPipe}'s reported ``$+1$ activation copy''
\citep[\S 3.2.1]{deepseek-v3} is a small price for halving the bubble.
The second is that \emph{\glsterm{fsdp}{FSDP} and \glsterm{tensor-parallel}{TP} both shard \glsterm{parameters}{parameters} but at
different granularities}: \glsterm{fsdp}{FSDP} collects per-step (one all-gather per
\texttt{DTensor} placement, amortized across the layer's compute), \glsterm{tensor-parallel}{TP}
collects per-layer (an all-reduce inside every linear). At small model
size \glsterm{fsdp}{FSDP}'s per-step collective dominates; at large model size \glsterm{tensor-parallel}{TP}'s
per-layer collective dominates --- which is why pure-\glsterm{fsdp}{FSDP} plans win on
sub-$10\text{B}$ models and \glsterm{tensor-parallel}{TP}+\glsterm{pipeline-parallel}{PP} plans dominate at $400\text{B}{+}$
\citep[\S 1; \S 2.1.3]{torchtitan2024}. \glsterm{expert-parallel}{EP} adds a third granularity:
the all-to-all is per-MoE-layer, so \glsterm{expert-parallel}{EP}-heavy plans on $L$-layer models
with $L_{\text{MoE}} \approx L$ pay $2 L$ all-to-alls per step
\citep[\S 3.2]{deepseek-v3}.

\begin{analogy}
The five axes lay a 5D crystal over the worker pool: each GPU sits at
one cell of a $d \times t \times c \times p \times e$ mesh, and the
``shape'' of the crystal --- which axes are nontrivial and how big they
are --- determines which collectives fire on which axis at which
frequency. Returning to canonical form: this is exactly the abstraction
\texttt{torch.distributed.DeviceMesh} represents \citep[\S
2.1.1]{torchtitan2024}; the five axes are quite literally the five
dimensions of a \texttt{DTensor} placement, and the analogy is
operational, not just suggestive. The parallelism plan's tuple
$(d, t, c, p, e)$ is the shape of the mesh; the placements declared on
each \texttt{Linear} / norm / attention module are the per-axis
collectives the runtime issues; the cost in \cref{tab:dt-memory} is the
amortized communication that mesh shape implies.
\end{analogy}

\subsection{Composability --- TorchTitan's 4D + Float8 + AsyncTP}
\label{sec:dt-torchtitan}

\citet{torchtitan2024} package the four non-\glsterm{expert-parallel}{EP} axes (\glsterm{data-parallel}{DP} / \glsterm{tensor-parallel}{TP}+\glsterm{sequence-parallel}{SP} / \glsterm{context-parallel}{CP} /
\glsterm{pipeline-parallel}{PP}) on a single 4D \texttt{DeviceMesh} with \glsterm{fsdp2}{FSDP2} as the default
1D path. The headline composability claim is that any subset of axes
can be enabled without retouching the model definition: parallelism is
declared on the mesh, not in the layers, and the runtime issues the
right collectives. The reported speedups, on optimized
\texttt{torch.compile}-based BF16 baselines on H100, are
``\textbf{$65.08\%$ on Llama 3.1 8B at 128 GPU scale (1D)},
\textbf{$12.59\%$ on Llama 3.1 70B at 256 GPU scale (2D)}, to
\textbf{$30\%$ on Llama 3.1 405B at 512 GPU scale (3D)}''
\citep[abstract]{torchtitan2024}.\footnote{KB excerpt:
\texttt{kb/excerpts/torchtitan2024.md\#abstract}.} The components that
ride on top of the mesh are themselves composable:
\emph{torchao.float8} for selective per-tensor \glsterm{fp8}{FP8} \texttt{Linear}
layers (\cref{sec:mixed-precision}), \emph{Async \glsterm{tensor-parallel}{TP}} for matmul-
collective overlap via \texttt{SymmetricMemory} intra-node buffers,
\emph{\glsterm{zerobubble}{ZeroBubble} / Flexible-Interleaved-1F1B} pipeline schedules
expressed as a list of compute actions plus compiler passes that insert
the corresponding collectives \citep[\S 2.1.4]{torchtitan2024}.\footnote{KB
excerpt: \texttt{kb/excerpts/torchtitan2024.md\#sec-2-1-4-pp}.}
\glsterm{fsdp2}{FSDP2}'s per-parameter \texttt{DTensor} sharding is the technology that
makes the rest compose: each parameter knows its own placement on the
mesh, so a \glsterm{tensor-parallel}{TP}-sharded \texttt{Linear} composes with \glsterm{fsdp2}{FSDP2} by stacking
placements rather than reconciling buffer layouts. \deepencite{FSDP2-
DTensor formal placement composition --- KB excerpt covers the
$7\%$/$1.5\%$ improvements but not the placement algebra in detail.}

\subsection{Case study --- DeepSeek-V3's plan}
\label{sec:dt-deepseek}

\glsterm{precision-pinning}{DeepSeek-V3} \citep{deepseek-v3} is the most documented frontier-scale
MoE training run as of 2026 --- $671\text{B}$ total / $37\text{B}$
active \glsterm{parameters}{parameters}, $14.8\text{T}$ tokens, $2{,}048$ H800 GPUs, and the
\emph{full} pre-training, context-extension, and post-training cost
table reproduced as \cref{tab:dsv3-cost}.

\begin{table}[h]
\centering
\small
\begin{tabular}{lrr}
\toprule
Stage & H800 GPU-hours & USD ($\$2$/GPU-hour) \\
\midrule
Pre-training         & $2{,}664$K & $\$5.328$M \\
\glsterm{context-extension}{Context extension}    & $\;\;119$K & $\$0.238$M \\
Post-training        & $\;\;\;\;5$K & $\$0.010$M \\
\midrule
\textbf{Total}       & $\mathbf{2{,}788}$\textbf{K} & $\mathbf{\$5.576}$\textbf{M} \\
\bottomrule
\end{tabular}
\caption{DeepSeek-V3 full training cost. ``Training DeepSeek-V3 on
each trillion tokens requires only 180K H800 GPU hours, i.e., 3.7
days on our cluster with 2048 H800 GPUs.'' After
\citet[\S 1]{deepseek-v3}.\footnotemark}
\label{tab:dsv3-cost}
\end{table}\footnotetext{KB excerpt:
\texttt{kb/excerpts/\glsterm{precision-pinning}{deepseek-v3}-training.md\#sec-1-cost}.}

The parallelism plan is \citep[\S 3.2]{deepseek-v3}\footnote{KB
excerpt:
\texttt{kb/excerpts/\glsterm{precision-pinning}{deepseek-v3}-training.md\#sec-3-2-parallelism}.}:
\textbf{16-way \glsterm{pipeline-parallel}{PP} with the \glsterm{dualpipe}{DualPipe} schedule, 64-way \glsterm{expert-parallel}{EP} spanning 8
nodes, ZeRO-1 \glsterm{data-parallel}{DP}, and \emph{no \glsterm{tensor-parallel}{TP}}.} Concretely, with $G = 2{,}048$,
$p = 16$, $e = 64$, no \glsterm{tensor-parallel}{TP} and no \glsterm{context-parallel}{CP}, $d = G / (p \cdot e) = 2$ at the
\glsterm{data-parallel}{DP} axis.\footnote{The plan is best read as: \glsterm{pipeline-parallel}{PP} $\times$ \glsterm{expert-parallel}{EP} carves the
$2{,}048$ workers into $2$ \glsterm{data-parallel}{DP} replicas of a $16 \times 64$ parallelism
plane; ZeRO-1 splits optimizer state across the 2-way \glsterm{data-parallel}{DP}. Each MoE
layer's experts are partitioned across the 64-way \glsterm{expert-parallel}{EP} plane; non-MoE
weights replicate within \glsterm{expert-parallel}{EP} and shard across \glsterm{pipeline-parallel}{PP}.} The deliberate
absence of \glsterm{tensor-parallel}{TP} is the load-bearing choice: ``we meticulously optimize
the memory footprint during training, thereby enabling us to train
\glsterm{precision-pinning}{DeepSeek-V3} \emph{without using costly Tensor Parallelism}''
\citep[\S 3.2]{deepseek-v3}. The activation relief that \glsterm{tensor-parallel}{TP} would
normally provide comes instead from (i) \glsterm{fp8}{FP8} forward activations
(\cref{sec:mixed-precision}), (ii) the \glsterm{dualpipe}{DualPipe} overlap between
MoE all-to-alls and \glsterm{pipeline-parallel}{PP} forwards/backwards, and (iii) \glsterm{precision-pinning}{DeepSeek-V3}'s
fine-grained MoE architecture (256 routed experts) that keeps each
expert's per-token compute small enough for one GPU.

\paragraph{DualPipe --- the bidirectional bubble-eater.} Each forward
+ backward chunk is split into four sub-stages: \texttt{attention},
\texttt{all-to-all dispatch}, \texttt{MLP}, \texttt{all-to-all
combine} \citep[\S 3.2.1]{deepseek-v3}. For a backward chunk, both
\texttt{attention} and \texttt{MLP} are further split into ``backward
for input'' and ``backward for weights'' (the \glsterm{zerobubble}{ZeroBubble} decomposition
that isolates $W$). The \glsterm{key}{key} \glsterm{dualpipe}{DualPipe} move is to schedule these
sub-stages so the four communication phases (the all-to-all dispatch +
combine on each microbatch's forward and backward) overlap with the
matrix-compute phases of \emph{adjacent microbatches}, and to feed the
pipeline from both ends simultaneously \citep[\S 3.2.1]{deepseek-v3}.
The bubble formula in \cref{eq:bubble-formulas} drops to
$(\tfrac{p}{2} - 1)(F\&B + B - 3W)$ --- roughly half of 1F1B at the
cost of ``$2\times$ parameter memory and $+1$ activation copy''
\citep[\S 3.2.1]{deepseek-v3}.

\paragraph{Cross-node all-to-all --- exploiting the $3.2\times$
NVLink/IB ratio.} Within a node, $8$ H800s are connected via NVLink
with $160$ GB/s; across nodes the connection is InfiniBand with
$50$ GB/s --- a $3.2\times$ ratio
\citep[\S 3.2.2]{deepseek-v3}.\footnote{KB excerpt:
\texttt{kb/excerpts/\glsterm{precision-pinning}{deepseek-v3}-training.md\#sec-3-2-2-allreduce}.} A
naive \glsterm{expert-parallel}{EP} all-to-all over $64$ workers spanning $8$ nodes would put
every token on the IB hop $7$ times. \glsterm{precision-pinning}{DeepSeek-V3}'s kernel rewrites the
all-to-all as a two-phase routing constraint:
\begin{equation}
\text{(token, dest node)} \;\to\; \underbrace{1\;\text{IB transmit}}_{\text{to GPU with matching in-node index}}
\;\to\; \underbrace{\text{NVLink forward}}_{\text{within destination node, to expert host}},
\label{eq:dsv3-routing}
\end{equation}
with the routing capped at \textbf{at most 4 destination nodes per
token}. Reading: each token pays one IB transmit per (token, target
node), then everything else is intra-node NVLink. Combined with
\glsterm{dualpipe}{DualPipe} overlap of the dispatch and combine phases against compute,
this lets each token select an average of $3.2$ experts per node ``at
no additional cost from NVLink'' \citep[\S 3.2.2]{deepseek-v3}, while
keeping IB the only inter-node bottleneck.

\begin{intuition}
The \glsterm{dualpipe}{DualPipe} schedule plus the IB+NVLink kernel together turn
\glsterm{precision-pinning}{DeepSeek-V3}'s nominally $1{:}1$ compute-to-communication ratio
\citep[\S 3.2.1]{deepseek-v3} into an effectively compute-bound run
with no irrecoverable loss spikes \citep[abstract]{deepseek-v3}.
Returning to canonical form: in \cref{eq:bubble-formulas}, $F\&B$ is
the overlapped forward-and-backward sub-step where the all-to-all
phases are scheduled into the gaps between matmul phases of adjacent
microbatches; in \cref{eq:dsv3-routing} the $4$-node cap turns each
token's $7$-hop IB cost into a $1$-hop IB cost plus intra-node NVLink
forwarding, exploiting exactly the $3.2\times$ bandwidth ratio. The
two optimizations work together: \glsterm{dualpipe}{DualPipe} gives the kernel a fixed
schedule of when each token's dispatch/combine fires, and the
$4$-node-cap rule gives the kernel a topology-aware route.
\end{intuition}

\subsection{Distributed Muon}
\label{sec:dt-muon}

The 2025 add to the distributed-training surface is \emph{Distributed
\glsterm{per-shape-rms-scaling}{Muon}} \citep[\S 2.3]{muon-moonlight2025}\footnote{KB excerpt:
\texttt{kb/excerpts/\glsterm{muon}{muon}-moonlight2025.md\#sec-2-3-distributed}.}: a
sharded variant of the \glsterm{per-shape-rms-scaling}{Muon} optimizer (\cref{sec:optimization})
compatible with ZeRO-1 \glsterm{data-parallel}{DP}. The wrinkle is that \glsterm{per-shape-rms-scaling}{Muon}'s Newton-Schulz
orthogonalization step (Eq.~(2) of \cref{sec:optimization}) needs the
\emph{full} matrix gradient $\mathbf{M}_t \in \R^{A \times B}$, while
ZeRO-1 hands each rank only its
$1/d$-th partition of the master weight and the corresponding
partition of the gradient. \citet[\S 2.3]{muon-moonlight2025} resolve
this with two added operations layered on top of the vanilla
ZeRO-1-\glsterm{adamw}{AdamW} pattern:
\begin{enumerate}
\item \textbf{\glsterm{data-parallel}{DP} gather.} For each local \glsterm{data-parallel}{DP}-partitioned master weight
($1/d$ the size of the model weight), gather the corresponding
partitioned gradients across all \glsterm{data-parallel}{DP} ranks into the full gradient matrix
$\mathbf{M}_t$.
\item \textbf{Calculate full update.} Run the \glsterm{newton-schulz-iteration}{Newton-Schulz iteration}
on $\mathbf{M}_t$, producing $\mathbf{O}_t \in \R^{A \times B}$;
\emph{discard} the part of $\mathbf{O}_t$ outside this rank's parameter
partition; apply only the local-partition update.
\end{enumerate}
The reported overhead is $1\!-\!3\%$ of step latency
\citep[\S 2.3]{muon-moonlight2025}, the only additional bandwidth cost
is the gradient gather, and the orthogonalization itself is data-
parallel: each rank does its own Newton-Schulz on the full matrix and
keeps only its slice. The Moonlight model (3B-active / 16B-total MoE,
$5.7\text{T}$ tokens) is the published demonstration that this
composes with the rest of the stack at frontier-MoE scale; whether it
holds at $70\text{B}{+}$ dense or $600\text{B}{+}$ MoE is part of the
optimizer-frontier open question discussed in \cref{sec:optimization}.

\subsection{Variants and lineage}
\label{sec:dt-lineage}

\cref{tab:dt-lineage} traces the system lineage. The 2020-era systems
forced one axis: Megatron-LM was \glsterm{tensor-parallel}{TP}+\glsterm{pipeline-parallel}{PP} (no \glsterm{fsdp}{FSDP}), DeepSpeed was
ZeRO-1/2/3 (no model parallel). The 2024-era systems --- TorchTitan
\citep{torchtitan2024} on the open side, the DeepSeek stack
\citep{deepseek-v3} on the case-study side --- treat each axis as an
independent dimension on a \texttt{DTensor} mesh and let users compose
them.

\begin{table}[h]
\centering
\small
\begin{tabular}{llll}
\toprule
System & Year & Distinctive choice & Anchor \\
\midrule
Megatron-LM (NVIDIA) & 2019-- & First \glsterm{tensor-parallel}{TP}+\glsterm{pipeline-parallel}{PP} at scale & \citep{shoeybi2019-megatron} \\
ZeRO-1/2/3 (DeepSpeed) & 2020--21 & Optim/grad/param-shard ladder & \citep{rajbhandari2020-zero} \\
FSDP1 (\glsterm{dtensor-devicemesh}{PyTorch}) & 2022 & Native ZeRO-3 with FlatParameter & \citep{zhao2023-fsdp} \\
Megatron-3D & 2021 & \glsterm{tensor-parallel}{TP} $\times$ \glsterm{pipeline-parallel}{PP} $\times$ \glsterm{data-parallel}{DP} composition & \citep{narayanan2021-megatron-3d} \\
GSPMD (XLA) & 2021 & Sharding propagation in compiler & \deepencite{Xu et al.\ 2021 GSPMD --- not in papers.json} \\
TorchTitan (\glsterm{dtensor-devicemesh}{PyTorch}) & 2024 & 4D + Float8 + \glsterm{asynctp}{AsyncTP} composable & \citep{torchtitan2024} \\
\glsterm{fsdp2}{FSDP2} (\glsterm{dtensor-devicemesh}{PyTorch}) & 2024 & Per-parameter DTensor sharding & \citep[\S 2.1.2]{torchtitan2024} \\
DeepSpeed-MoE / Tutel & 2022--23 & \glsterm{expert-parallel}{EP} + all-to-all dispatch kernels & \deepencite{Rajbhandari et al.\ 2022 DeepSpeed-MoE / Hwang et al.\ 2023 Tutel --- not in papers.json} \\
\glsterm{precision-pinning}{DeepSeek-V3} / \glsterm{dualpipe}{DualPipe} & 2024 & Bidirectional \glsterm{pipeline-parallel}{PP} + ZB decomposition & \citep[\S 3.2.1]{deepseek-v3} \\
Distributed \glsterm{per-shape-rms-scaling}{Muon} (Moonlight) & 2025 & Per-rank Newton-Schulz on gathered grad & \citep[\S 2.3]{muon-moonlight2025} \\
\bottomrule
\end{tabular}
\caption{Distributed-training systems lineage. The arc is toward
\emph{composability}: every axis available, declared via mesh
placement rather than wired into the model definition.}
\label{tab:dt-lineage}
\end{table}

The auto-parallelism direction --- pick $(d, t, c, p, e)$ from
$(N, D, \text{cluster shape})$ from the model graph --- is the open
research surface that GSPMD and Alpa moved into.
\deepencite{Zheng et al.\ 2022 ``Alpa: Automating Inter- and
Intra-Operator Parallelism for Distributed Deep Learning'' --- KB
note refs but excerpt-key not yet in papers.json.} Neither has become
the dominant tool of 2026 open practice; frontier plans are still
hand-tuned by ML systems engineers.

\subsection{Three live contradictions (set up; full discussion §14)}
\label{sec:dt-contradictions}

\begin{contradiction}
\textbf{\glsterm{tensor-parallel}{TP}-yes-vs-\glsterm{tensor-parallel}{TP}-no.} \citet{meta-llama3} train Llama 3 405B with
heavy \glsterm{tensor-parallel}{TP}+\glsterm{pipeline-parallel}{PP}+\glsterm{data-parallel}{DP} --- the prototypical Megatron-3D shape. \citet{deepseek-v3}
train \glsterm{precision-pinning}{DeepSeek-V3} 671B-A37B \emph{without \glsterm{tensor-parallel}{TP} at all}, claiming the
combination of ZeRO-1 \glsterm{data-parallel}{DP} + \glsterm{expert-parallel}{EP} + \glsterm{dualpipe}{DualPipe} + \glsterm{fp8}{FP8} supplies enough
activation relief to make \glsterm{tensor-parallel}{TP}'s per-linear collective cost a net loss.
This is a real architectural-conditional disagreement, not a wash. \glsterm{tensor-parallel}{TP}
costs $2L$ inter-GPU all-reduces per step; \glsterm{expert-parallel}{EP} costs $2L_{\text{MoE}}$
inter-node all-to-alls per step. DeepSeek's fine-grained MoE
architecture forces them to pay the \glsterm{expert-parallel}{EP} all-to-all already
\citep[\S 3.2]{deepseek-v3}, and pairing it with \glsterm{dualpipe}{DualPipe} overlap +
the $4$-node-cap kernel hides the cost. Llama 3 is dense at active-
parameter level; without an \glsterm{expert-parallel}{EP} all-to-all to fold \glsterm{tensor-parallel}{TP} work into, the
Megatron-3D plan stays preferable. The general rule is that
\emph{which} parallelism is best is a function of the model graph
(MoE vs.\ dense, expert count, expert size, attention head count) ---
not a universal answer.
\end{contradiction}

\begin{contradiction}
\textbf{Optimal plan from the model graph.} No public tool reliably
picks $(d, t, c, p, e)$ from $(N, D, \text{cluster shape}, \text{model
graph})$. GSPMD propagates user-declared sharding annotations through
the XLA computation graph, but it does not search; AlpaDiscovery
attempts joint inter- and intra-operator search but is not the
dominant tool in 2026 open practice. \citet[\S 1]{torchtitan2024}
explicitly punts: ``\emph{combining different parallelism methods
manually [is] complex and error-prone.}'' Frontier plans are still
hand-tuned by ML systems engineers iterating on memory traces and \glsterm{mfu}{MFU}
profiles. A theoretically grounded plan-from-graph optimizer is open.
\end{contradiction}

\begin{contradiction}
\textbf{Blackwell NVL72 inversion.} The \glsterm{precision-pinning}{DeepSeek-V3} cross-node kernel
\citep[\S 3.2.2]{deepseek-v3} exploits a specific topology fact: 8 GPUs
per node on NVLink ($160$ GB/s intra-node), nodes connected by IB
($50$ GB/s inter-node). The ratio $160/50 \approx 3.2\times$ motivates
the at-most-4-nodes-per-token routing cap, which keeps IB the binding
constraint and amortizes intra-node NVLink. NVIDIA's GB200 NVL72 rack
puts \emph{72 GPUs} in a single NVLink domain via NVSwitch, with
nominal $1{,}800$ GB/s effective per-GPU NVLink bandwidth across the
full $72$-GPU mesh. The IB-vs-NVLink ratio that DeepSeek's kernel
exploits ($3.2\times$ within an 8-GPU NVLink island) inverts: at NVL72
scale, an entire 72-GPU rack is one NVLink domain, only cross-rack
hops cross IB, and the optimal kernel is no longer ``cap each token
to 4 nodes'' but ``saturate the 72-GPU NVLink domain and minimize
cross-rack hops.'' The 2026-frontier hardware shift will require new
topology-aware all-to-all kernels; the current generation's are not a
universal solution. \deepencite{NVL72 effective NVLink bandwidth +
re-derivation of token-routing cost under one-rack-NVLink topology
--- not yet in papers.json; tracked via NVIDIA Blackwell whitepapers.}
\end{contradiction}

\subsection{Forward link}
\label{sec:dt-forward}

This section's parallelism plan, the optimizer choice from
\cref{sec:optimization}, and the data and scaling-law choices from
\cref{sec:pre-training-data,sec:scaling-laws} together fix
\emph{everything except the precision axis} of the recipe. The next
section (\cref{sec:mixed-precision}) closes the recipe by walking
that precision axis: storage, compute, communication, and accumulator
precision as four independently tunable choices; the FP32-master-plus-
low-precision-working pattern \citep[\S 2]{micikevicius2017-mixed-precision};
BF16 as a drop-in replacement for FP32 that abolishes \glsterm{loss-scaling}{loss scaling}
\citep{kalamkar2019-bfloat16}; and \glsterm{precision-pinning}{DeepSeek-V3}'s fine-grained \glsterm{fp8}{FP8}
recipe \citep[\S 3.3]{deepseek-v3} that --- combined with the
parallelism plan above --- delivered $14.8\text{T}$ tokens of pre-
training with zero irrecoverable loss spikes
\citep[abstract]{deepseek-v3}. The connection between this section and
the next is structural: TorchTitan's per-tensor Float8 path
\citep[\S 2.2.4]{torchtitan2024} composes against the 4D
\texttt{DeviceMesh}, \glsterm{precision-pinning}{DeepSeek-V3}'s per-tile \glsterm{fp8}{FP8} recipe composes
against the cross-node all-to-all kernel from \cref{sec:dt-deepseek},
and the optimizer-precision coupling first surfaces here (Distributed
\glsterm{per-shape-rms-scaling}{Muon}'s full-matrix $\mathbf{O}_t$ exceeds BF16 high-precision range
without weight decay; \cref{sec:mixed-precision} works it through to
the BF16 / \glsterm{fp8}{FP8} recipe).
